\section{IA -- Recherche taboue}

\subsection{Adaptation de l'algorithme à l'ordonnancement}

\aremplir{définir la représentation d'une solution, le voisinage, la liste taboue, le critère d'aspiration et les conditions d'arrêt.}

\subsection{Base de données de petite taille}

\subsubsection{Courbes et tableaux}
\aremplir{présenter les résultats reproductibles.}

\subsubsection{Analyse des résultats}
\aremplir{interpréter la qualité, la convergence et le temps de calcul.}

\subsubsection{Intérêt de l'étude}
\aremplir{expliquer ce que l'instance contrôlée permet de valider.}

\subsection{Base de données de grande taille}

\subsubsection{Courbes et tableaux}
\aremplir{présenter les résultats reproductibles.}

\subsubsection{Analyse des résultats}
\aremplir{interpréter la qualité, la robustesse et le passage à l'échelle.}

\subsubsection{Intérêt de l'étude}
\aremplir{relier les résultats aux contraintes opérationnelles.}
